% 第18批：学习指导自测题7选择题7.1.1--7.1.8；源印刷167--168／PDF175--176。
\subsection{自测题7：静电、电容、载流导线和位移电流选择题}

\sourcetest{7.1.1}{167}{175}
\problemtext 两块面积均为 $S$、相距 $d$ 的大型平行金属板，带总电量 $q_1,q_2$；求两板电势差 $U_{AB}$。
\workings 忽略边缘效应。平行板间的合场强沿 $A\to B$ 的有向分量为
\[E_x=\frac{q_1-q_2}{2\varepsilon_0 S}.\]
所以
\[\boxed{U_{AB}=V_A-V_B=\frac{(q_1-q_2)d}{2\varepsilon_0 S}}.\]
按原选项为 \textbf{C}。
\crosscheck 不能将两板的总电荷都当成同一侧的高斯面电荷，亦不能将电势差错误写成与 $q_1+q_2$ 成正比。

\sourcetest{7.1.2}{167}{175}
\problemtext 题7.1.2示意均匀正带电圆盘的电场线；比较圆盘边缘外 $P$ 与圆盘中心 $O$ 的电势。
\workings 静电场线从正电荷出发并指向低电势方向。图示由 $O$ 附近向外经过 $P$ 的场线表明
\[\boxed{V_P<V_O}.\]
按选项为 \textbf{B}。
\crosscheck 只需沿电场线的电势下降性质，不需要先知道圆盘轴外任意点的精确场强解析式。

\sourcetest{7.1.3}{167}{175}
\problemtext 空气平行板电容器两极板面积均为 $S$，分别带 $\pm q$；忽略边缘效应，求极板间吸引力。
\workings 一个极板仅受\emph{另一}极板所产生的场 $E_{\rm other}=\sigma/(2\varepsilon_0)$。故
\[\boxed{F=qE_{\rm other}=\frac{q^2}{2\varepsilon_0 S}},\qquad \sigma=q/S.\]
选 \textbf{B}。
\crosscheck 电容器板间总场为 $\sigma/\varepsilon_0$，但不能把板自身场也计入自身受力，否则多出因子2。

\sourcetest{7.1.4}{167}{175}
\problemtext 相同电流 $I$ 分别流过直径为 $a$ 的圆线圈与边长为 $a$ 的正方形线圈，求中心磁场大小比 $B_{\rm circ}/B_{\rm square}$。
\workings 单匝圆线圈中心 $B_{\rm circ}=\mu_0I/a$。正方形的每条边到中心垂距为 $a/2$，两端夹角均为 $45^\circ$；按有限直导线公式四边叠加，
\[B_{\rm square}=\frac{2\sqrt2\,\mu_0I}{\pi a}.\]
因此
\[\boxed{\frac{B_{\rm circ}}{B_{\rm square}}=\frac{\pi}{2\sqrt2}\simeq1.11}.\]
选 \textbf{C}。
\crosscheck 题设圆的\emph{直径}等于方形\emph{边长}，不是圆半径等于边长。

\sourcetest{7.1.5}{167}{175}
\problemtext 题7.1.3图中四个相同质量、等电荷绝对值的粒子从 $O$ 点以不同速率射入匀强磁场，$B$ 垂直纸面向外；辨认动能最大的负电粒子轨迹。
\workings 同质量、同 $|q|$ 时，轨道曲率半径为
\[r=\frac{mv}{|q|B},\qquad K=\frac{q^2 B^2r^2}{2m}.\]
粒子初速度向左，带负电者受力初始指向纸面下方，对应图中 $Oc,Od$；其中 $Od$ 偏折较小，半径更大，因此动能最大者为 \textbf{轨迹 $Od$（D）}。
\crosscheck 轨道向上/向下辨认电荷正负，曲率半径辨认动能；单凭轨迹长度不可判定动能。

\sourcetest{7.1.6}{168}{176}
\problemtext 题7.1.4图：固定载流大平板附近放置一小线框，线框平面与平板垂直；判断框受磁场作用后的运动趋势。
\workings 大平板近似匀强磁场时，闭合小线圈的磁偶极矩为 $\bm m=I_2S\,\hat{\bm n}$，受力矩
\[\boxed{\bm \tau=\bm m\times\bm B}.\]
磁场近似空间均匀时总平动力为零，而力矩使磁矩趋向沿磁场方向排列，故\emph{应以转动为主而非平移}。源图绕行箭头方向的透视存在辨识风险，\textbf{顺时针/逆时针选项暂不定稿}。
\crosscheck 待放大原图确定大平板正面、线框电流环向和从板向外的观察方向，再确定转动的有向手性；不能只由受力矩公式猜 A/B/C/D。

\sourcetest{7.1.7}{168}{176}
\problemtext 理想环形螺线管电感系数为 $L$，从对称位置分切为两个半环，各自电感与 $L$ 的关系。
\workings 取截面积、磁导率不变，理想环形线圈自感 $L\propto N^2/\ell$。各半环有 $N/2$ 匝、磁路长度 $\ell/2$，
\[\boxed{L_{\rm half}=L\,\frac{(1/2)^2}{1/2}=\frac L2}.\]
故两半环各为 $L/2$，选 \textbf{A}。
\crosscheck 半环的匝数减半同时磁路长度也减半；单纯按 $N^2$ 缩放为 $L/4$ 忽略磁路长度变化。开口边缘漏磁不属于理想近似。

\sourcetest{7.1.8}{168}{176}
\problemtext 判断关于位移电流的起因、发热效应与安培环路定理的四种表述。
\workings 麦克斯韦位移电流密度定义为 $\bm J_d=\partial\bm D/\partial t$，由变化的电位移通量产生，故对应“变化电场产生位移电流”的 \textbf{A}。
\[\boxed{I_d=\frac{d\Phi_D}{dt}}.\]
\crosscheck 位移电流不是材料中载流子定向运动，不能照搬导电电流的焦耳热公式；安培--麦克斯韦环路定律明确包含位移电流。
